%%%%%%%%%%%%%%%%%%%%%%
% This file is a self-contained blueprint for the Gowers-Hatami theorem, following Thomas' lecture notes.
% Many definitions have been formalized in representation theory in mathlib:
% https://leanprover-community.github.io/mathlib_docs/representation_theory/basic.html
% However for completeness we include the definitions as comments.
%%%%%%%%%%%%%%%%%%%%%%

\section{Preliminaries}
\begin{definition}[$\sigma$ inner product.]\label{sigmaInner}\lean{sigmaInner}\leanok
    The $\sigma$ inner product between matrices $A$ and $B$ is defined as
    \be
    \inner{A}{B}_{\sigma} = \Tr(A^* B \sigma),
    \ee
    where $\sigma \geq 0$ and $A^*$ denotes conjugate-transpose.
\end{definition}

\begin{definition}[$(\epsilon, \sigma)$-representation.]\label{approx-representation}\uses{sigmaInner}\lean{IsApproxRepresentation}\leanok
Given a finite group \(G\), an integer \(d \geq 1, \varepsilon \geq 0\), and a \(d\)-dimensional positive semidefinite matrix \(\sigma\) with trace 1 , an \((\varepsilon, \sigma)\)-representation of \(G\) is a map \(f: G \rightarrow U_{d}(\mathbb{C})\), to the  \(d \times d\) unitary matrices, such that
\be
\mathbb{E}_{x, y \in G} \Re\left(\left\langle f(x)^{*} f(y), f\left(x^{-1} y\right)\right\rangle_{\sigma}\right) \geq 1-\varepsilon,
\ee
where $\Re$ denotes taking real part of the inner product. 
\end{definition}

\begin{proposition}[$\sigma$-norm Square]\label{sigmaNormSqSub}\uses{sigmaInner}\lean{sigmaNormSq_sub_eq}\leanok
Let $\sigma$ be positive semidefinite and let $A, B$ be matrices of the same dimension. Then
\[
\|A - B\|_\sigma^2 = \|A\|_\sigma^2 + \|B\|_\sigma^2 - 2\,\Re\inner{A}{B}_\sigma.
\]
\end{proposition}
\begin{proof}\leanok
    Expanding by linearity of the inner product,
    \[
    \|A-B\|_\sigma^2 = \Re\inner{A-B}{A-B}_\sigma
    = \|A\|_\sigma^2 + \|B\|_\sigma^2 - \Re\inner{A}{B}_\sigma - \Re\inner{B}{A}_\sigma.
    \]
    Since $\sigma$ is Hermitian, $\inner{B}{A}_\sigma = \overline{\inner{A}{B}_\sigma}$, so the two cross terms have the same real part, yielding $-2\,\Re\inner{A}{B}_\sigma$.
\end{proof}


\section{Representation theory}

% Representation theory exists on mathlib but for completeness we include the definitions as comments

% \begin{definition}[Representation]
%     Let $G$ be a group and $V$ a finite-dimensional complex vector space. A \emph{representation} of $G$ is a homomorphism
%     \be
%         \rho : G \to \mathrm{GL}(V),
%     \ee
%     such that for all $x,y \in G$,
%     \be
%         \rho(xy) = \rho(x)\rho(y).
%     \ee
% \end{definition}

In general, a representation of a group \(G\) on a vector space \(V\) over a
field \(k\) is a homomorphism from \(G\) to the linear endomorphisms of \(V\).
In mathlib this is the abbreviation \(\texttt{Representation}\), defined as
\(G \to^* (V \to_{\!k} V)\), i.e. a monoid homomorphism into the \(k\)-linear
endomorphisms of \(V\).

We call a representation \emph{unitary} when its image lies in the unitary
group; in the finite-dimensional matrix setting used below this means a
homomorphism \(\rho : G \to U(V)\). In the Lean development this is represented
directly as a monoid homomorphism into a matrix unitary group, for example
\(G \to^* \texttt{Matrix.unitaryGroup}\).

% \begin{definition}[Irreducible representation]
%     A representation $\rho : G \to \mathrm{GL}(V)$ is called \emph{irreducible} if the only subspaces $W \subseteq V$ satisfying
%     \be
%         \rho(x) W \subseteq W \quad \text{for all } x \in G
%     \ee
%     are $W = \{0\}$ and $W = V$.
% \end{definition}

% For the rest of the notes, we will write ``irrep'' as a shorthand to denote irreducible representation.

% \begin{definition}[Completely reducible]
%     A representation $\rho : G \to \mathrm{GL}(V)$ is \emph{completely reducible} if $V$ can be decomposed as a direct sum of invariant subspaces,
%     \be
%         V = \bigoplus_i V_i,
%     \ee
%     where each $V_i$ carries an irreducible representation of $G$.
% \end{definition}

% \begin{definition}[Inequivalent representations]
%     Two representations $\rho_1 : G \to \mathrm{GL}(V)$ and $\rho_2 : G \to \mathrm{GL}(W)$ are \emph{equivalent} if there exists an invertible linear map $T : V \to W$ such that
%     \be
%         T \rho_1(x)T^{-1} = \rho_2(x) \quad \text{for all } x \in G.
%     \ee
%     They are \emph{inequivalent} if no such map exists.
% \end{definition}

% \begin{definition}[Character]
%     Let $\rho : G \to \mathrm{GL}(V)$ be a representation. The \emph{character} of $\rho$ is the function $\chi_\rho : G \to \mathbb{C}$ defined by
%     \be
%         \chi_\rho(x) = \mathrm{Tr}[\rho(x)].
%     \ee
% \end{definition}

\begin{definition}[Regular representation]\label{RegularRep}
    Let $R: G \rightarrow \mathbb{C}[G]$ be a regular representation where $\{\ket{g}, g\in G\}$ forms a basis of \(\mathbb{C}[G]\). The representation is given by
    \be
    R(x) \ket{g} = \ket{xg}.
    \ee
    
\end{definition}

It is possible to decompose any representation of a finite group into direct sums over irreps. 

\begin{theorem}[Maschke's theorem]\label{Maschke}
    Let $G$ be a finite group and let $\rho : G \to \mathrm{GL}(V)$ be a finite-dimensional representation over $\mathbb{C}$. Then $\rho$ is completely reducible. In particular, there exists a decomposition
    \be
        \rho \cong \bigoplus_{m} r_m \, \rho_m,
    \ee
    where $\{\rho_m\}$ is a set of inequivalent irreducible representations of $G$, and $r_m \in \mathbb{Z}^+$ are the multiplicities.
\end{theorem}

\begin{sproof}
    Maschke's theorem is partially formalized in $\texttt{Mathlib.RepresentationTheory.Maschke}$ but incomplete.
    Filippo Berta, Jonathan Conrad, and Younous Malliefer are formalizing this.
\end{sproof}

% mathlib: `FDRep.char_orthonormal`
% \begin{theorem}[Petit orthogonality theorem]\label{PetitOrthogonality}
%      Let $R_a$ and $R_b$ be two non-equivalent irreps of $G$. The characters satisfy the orthogonality relation
% \be
% \sum_{g\in G} \chi^*_a(g) \chi_b(g) = \abs{G} \delta_{a,b}.
% \ee    
% \end{theorem}

\begin{proposition}[Decomposition of the regular representation]\label{RegDecompIrrep}\uses{Maschke,RegularRep}
    Let $R : G \to \mathrm{GL}(\mathbb{C}[G])$ be the regular representation of a finite group $G$. Then
    \be
        R \cong \bigoplus_{\rho \in \widehat{G}} d_\rho \, \rho,
    \ee
    where $\widehat{G}$ is a complete set of inequivalent irreducible representations of $G$, and $d_\rho = \dim(V_\rho)$ is the dimension of $\rho$.
\end{proposition}

\begin{sproof}
    This can be proven using Maschke's theorem in \ref{Maschke} to decompose into irreps with multiplicity. 
    Then use character orthogonality theorem in mathlib $\texttt{FDRep.char\_orthonormal}$ to show the multiplicity equals irrep dimension. 
    A more formal statement is given by the Peter-Weyl theorem. 
\end{sproof}


\begin{proposition}[Character of regular representation.]\label{regular_repr_character}\uses{RegDecompIrrep}
Let $G$ be a finite group and $\sum_\rho$ denotes summing over the complete set of inequivalent irreps $\hat{G}$
    \be\label{eq:regular_repr_characterr}
    \sum_{\rho \in \hat{G}} d_{\rho} \operatorname{Tr}(\rho(x))=|G| \delta_{x, e}.
    \ee
\end{proposition}
\begin{sproof}
    The proof uses decomposition of regular representation and definition of character. 
\end{sproof}

% \begin{fact}
%     Let $\hat{G}$ be the complete set of inequivalent irreps of $G$ and $d_\rho = \dim(\rho)$ the dimension, then
%      \be\label{eq:dim_rho_group}
%      \sum_{\rho} d_{\rho}^{2}=|G|.
%      \ee
% \end{fact}
% \begin{sproof}
%     This is a corollary from Fact.~\ref{regular_repr_character} for $x=e$.
% \end{sproof}

\begin{definition}[Group fourier transform]\label{GroupFourier}
    Let $f: G \rightarrow U_d(\cc)$ be a map and \(\rho: G \rightarrow U_{d_{\rho}}(\mathbb{C})\) be a unitary irrep. The fourier transform of $f$ at the representation $\rho$ is denoted by $\hat{f}$:
    \be
\hat{f}(\rho)=\mathbb{E}_{x \in G} f(x) \otimes \overline{\rho(x)},
    \ee
    where $\overline{\rho(x)}$ denotes complex conjugate of $\rho(x)$. 
\end{definition}

\section{Gowers Hatami Theorem}

\begin{theorem}[Gowers-Hatami]\label{gowers-hatami}
  \uses{sigmaNormSqSub, approx-representation,GroupFourier, regular_repr_character}
Let \(G\) be a finite group, \(\varepsilon \geq 0\), and \(f: G \rightarrow U_{d}(\mathbb{C})\) an \((\varepsilon, \sigma)\)-representation of \(G\). Then there exists a \(d^{\prime} \geq d\), an isometry \(V: \mathbb{C}^{d} \rightarrow \mathbb{C}^{d^{\prime}}\), and a representation \(g: G \rightarrow U_{d^{\prime}}(\mathbb{C})\) such that
\be
\mathbb{E}_{x \in G}\left\|f(x)-V^{*} g(x) V\right\|_{\sigma}^{2} \leq 2 \varepsilon.
\ee 
\end{theorem}

\begin{sproof}
    The theorem is due to Gowers and Hatami~\cite{GowHat17}; the constructive proof below follows Vidick's lecture notes, Theorem 2.3~\cite{VidickQMATH22}.
    In particular, the proof is constructive. The isometry is defined as 
    \(V: \mathbb{C}^{d} \rightarrow \mathbb{C}^{d} \otimes\left(\oplus_{\rho} \mathbb{C}^{d_{\rho}} \otimes \mathbb{C}^{d_{\rho}}\right)\) by
    \be\label{eq:isometry_V}
    V u = \bigoplus_{\rho} d_{\rho}^{1 / 2} \sum_{i=1}^{d_{\rho}}\left(\hat{f}(\rho)\left(u \otimes e_{i}\right)\right) \otimes e_{i}.
    \ee
    And the exact representation is defined as 
    \be
    g(x)=\bigoplus_{\rho}\left(I_{d} \otimes I_{d_{\rho}} \otimes \rho(x)\right).
    \ee

    The proof requires the following two ingredients:
    \begin{enumerate}
        \item $V$ is an isometry:
        \be
        V^{*} V  = \I_d
        \ee 
        \item The isometry ``averages'' over group multiplication. For any \(x \in G\),
        \be
        V^{*} g(x) V = \mathbb{E}_{z} f(z)^{*} f(z x).
        \ee
    \end{enumerate}
    Both ingredients follow from the key identity over the regular representation in \equref{eq:regular_repr_characterr}.
\end{sproof}

Next, we will give a different proof that does not directly decompose into irreps but instead uses the action of the regular representation.
\begin{lemma}[Gowers-Hatami Pair-Correlation]\label{gowers-hatami-correlation}\lean{gh2_average_correlation, gh2_compression_avg_eq_pair, gh2_pair_correlation}\leanok\uses{approx-representation}
Let \(G\) be a finite group, \(\varepsilon \geq 0\), and \(\rho: G \rightarrow U(\mathcal{H})\) an \((\varepsilon, \sigma)\)-representation of \(G\) on space $\mathcal{H}$. 
Define the pair-correlation of $\rho$ as:
\be
    \mathcal{C}(\rho) := \EE_{x,y \in G} \Re \langle \rho(y)\rho(x), \rho(yx) \rangle_\sigma.
\ee
Then there exists an isometry \(V: \mathcal{H} \rightarrow \mathcal{H}^\prime\) to a different space $\mathcal{H}^\prime$, 
and an exact representation \(\rho^\prime: G \rightarrow U(\mathcal{H}^\prime)\) such that the average correlation between $\rho$ and $\rho^\prime$ is exactly the pair-correlation, which is bounded by $1-\varepsilon$:
\be
    \EE_{x \in G} \Re \langle \rho(x), V^{*} \rho^\prime(x) V\rangle_{\sigma} = \mathcal{C}(\rho) \geq 1-\varepsilon.
\ee    
\end{lemma}

\begin{proof}\leanok
    Here $\mathcal{H}^\prime = L(G, \mathcal{H})$ is the space of functions $f: G \to \mathcal{H}$.
    
    Let the isometry be defined by the action of the approximate representation on a state $u\in \mathcal{H}$ as
    \bea
    &V:& \mathcal{H} \to L(G, \mathcal{H}), \\
    & V(u):& G \to \mathcal{H}, \quad x \mapsto \rho(x)(u), \quad \forall x \in G.
    \eea

    The adjoint map $V^*$ is defined as the unique map such that for all $u \in \mathcal{H}$ and $f \in L(G, \mathcal{H})$,
    \be
    \langle V(u), f \rangle_{L(G, \mathcal{H})} = \langle u, V^*(f) \rangle_{\mathcal{H}}.
    \ee

    In particular, $V^*$ is explicitly given by the group average:
    \bea
    &V^*:& L(G, \mathcal{H}) \to \mathcal{H}, \\
    &V^*(f)& = \EE_{x \in G} \left[ \rho^*(x) f(x) \right].
    \eea

    The exact representation is given by the right regular representation on $L(G, \mathcal{H})$:
    \bea
    &\rho^\prime:& G \to U(L(G, \mathcal{H})), \\
    &(\rho^\prime(x)f)(y)& = f(yx), \quad \forall x,y \in G.
    \eea

    $V$ is an isometry because $\rho(x)$ is unitary, yielding:
    \be
    V^* V(u) = \EE_{x \in G} \left[ \rho^*(x) \rho(x)(u) \right] = \EE_{x \in G} [u] = u.
    \ee

    The compression of the exact representation ``averages'' over group multiplication. For any $u \in \mathcal{H}$:
    \bea\label{eq:isometry_average} 
    V^* \rho^\prime(x) V(u) & = & \EE_{y \in G} \left[ \rho^*(y) (\rho^\prime(x)\rho(y))(u) \right] \nonumber \\
    & = & \EE_{y \in G} \left[ \rho^*(y) \rho(yx)(u) \right].
    \eea

    Substituting the average property from Eq.~\ref{eq:isometry_average} into the inner product, and using the unitarity of $\rho(y)$ to rewrite the adjoints ($\langle A, B^* C \rangle_\sigma = \langle B A, C \rangle_\sigma$), we obtain:
    \bea \label{eq:inner_product_bound}
    \EE_{x \in G} \Re \langle \rho(x), V^* \rho^\prime(x) V\rangle_\sigma 
    & = & \EE_{x,y \in G} \Re \langle \rho(x), \rho^*(y) \rho(yx) \rangle_\sigma \nonumber \\ 
    & = & \EE_{x,y \in G} \Re \langle \rho(y)\rho(x), \rho(yx) \rangle_\sigma \nonumber \\ 
    & = & \mathcal{C}(\rho).
    \eea    
    By the definition of an $(\varepsilon, \sigma)$-approximate representation, the pair-correlation inherently satisfies $\mathcal{C}(\rho) \geq 1 - \varepsilon$, which concludes the proof.
\end{proof}
% Below is older lemma that was refactored above
% \begin{lemma}[Gowers-Hatami-Correlation]\label{gowers-hatami-correlation}\lean{gh2_average_correlation}\leanok\uses{approx-representation}
% Let \(G\) be a finite group, \(\varepsilon \geq 0\), and \(\rho: G \rightarrow U(\mathcal{H})\) an \((\varepsilon, \sigma)\)-representation of \(G\) on space $\mathcal{H}$. 
% Then there exists an isometry \(V: \mathcal{H} \rightarrow \mathcal{H}^\prime\) to a different space $\mathcal{H}^\prime$, 
% and a representation \(\rho^\prime: G \rightarrow \mathcal{H}^\prime\) such that the average correlation between $\rho$ and $\rho^\prime$ is:
% \be
% \mathbb{E}_{x \in G} \Re \langle f(x), V^{*} \rho^\prime(x) V\rangle_{\sigma} \geq 1-\varepsilon,
% \ee    
% \end{lemma}

% \begin{proof}
%     Here $\mathcal{H}^\prime = L(G, \mathcal{H})$ is the space of functions $f: G \to \mathcal{H}$.
    
%     Let the isometry be defined by the action of the approximate representation on a state $u\in \mathcal{H}$ as
%     \bea
%     &V:& \mathcal{H} \to L(G, \mathcal{H}), \\
%     & V(u):& G \to \mathcal{H}, \quad x \mapsto \rho(x)(u), \forall x \in G.
%     \eea

%     The adjoint map $V^*$ is defined as the unique map such that for all $u \in \mathcal{H}$ and $g \in L(G, \mathcal{H})$,
%     \be
%     \langle V(u), g \rangle_{L(G, \mathcal{H})} = \langle u, V^*(g) \rangle_{\mathcal{H}}.
%     \ee

%     In particular, $V^*$ is explicitly given by the group average
%     \bea
%     &V^*:& L(G, \mathcal{H}) \to \mathcal{H}, \\
%     &V^*(f)& = \EE_{x \in G} \left[ \rho^*(x) f(x) \right].
%     \eea

%     The exact representation is given by the right regular representation on $L(G, \mathcal{H})$:
%     \bea
%     &\rho^\prime:& G \to \End(L(G, \mathcal{H})), \\
%     &(\rho^\prime(x)f)(y)& = f(yx), \quad \forall x,y \in G.
%     \eea

%     $V$ is an isometry because
%     \be
%     V^* V(u) = \EE_{x \in G} \left[ \rho^*(x) \rho(x)(u) \right] = u.
%     \ee

%     The isometry ``averages'' over group multiplication because
%     \bea\label{eq:isometry_average} 
%     V^* \rho^\prime(x) V(u) & = & \EE_{y \in G} \left[ \rho^*(y) (\rho^\prime(x)\rho(y))(u) \right] \\
%     & = & \EE_{y \in G} \left[ \rho^*(y) (\rho(yx)(u)) \right].\nonumber
%     \eea

%     Finally, using the isometry average property in Eq.~\ref{eq:isometry_average} and the definition of $(\varepsilon, \sigma)$-representation, we have
%     \bea \label{eq:inner_product_bound}
%     && \EE_{x \in G} \Re \langle \rho(x), V^* \rho^\prime(x) V\rangle_\sigma, \\ \nonumber
%     & = & \EE_{x,y \in G} \Re \langle \rho(x), \rho^*(y) \rho(yx) \rangle_\sigma, \\ \nonumber
%     &\ge & 1 - \varepsilon. 
%     \eea    
% \end{proof}

\begin{theorem}[Gowers-Hatami-2]\label{gowers-hatami-2}\lean{gowers_hatami_GH2_product}\leanok
  \uses{sigmaNormSqSub, approx-representation}
Let \(G\) be a finite group, \(\varepsilon \geq 0\), and \(\rho: G \rightarrow U(\mathcal{H})\) an \((\varepsilon, \sigma)\)-representation of \(G\) on space $\mathcal{H}=\mathbb{C}^{d}$. 
Then there exists an isometry \(V: \mathcal{H} \rightarrow \mathcal{H}^\prime\) to a different space $\mathcal{H}^\prime$, 
and a representation \(\rho^\prime: G \rightarrow \mathcal{H}^\prime\) such that
\be
\mathbb{E}_{x \in G}\left\|\rho(x)-V^{*} \rho^\prime(x) V\right\|_{\sigma}^{2} \leq 2 \varepsilon.
\ee 
\end{theorem}

\begin{proof}\leanok
    \uses{gowers-hatami-correlation}
    \leanok
    A regular-representation proof of the Gowers-Hatami theorem appears in Morel's slides~\cite{Morel20}, following the original Gowers-Hatami theorem~\cite{GowHat17}.

    The rest of the proof follows by expanding the $\sigma$-norm, bounding the norm of $\EE_{x \in G} V^* \rho^\prime(x) V$ and applying the definition of $(\varepsilon, \sigma)$-representation.
     In particular, we have 
     \bea
        & &\EE_{x \in G}\left\|\rho(x)-V^{*} \rho^\prime(x) V\right\|_{\sigma}^{2} \\ 
        &= & \EE_{x \in G} \| \rho(x) \|_\sigma^2 + \EE_{x \in G} \|V^* \rho^\prime(x) V\|_\sigma^2 - 2 \EE_{x \in G} \Re \langle \rho(x), V^* \rho^\prime(x) V\rangle_\sigma \nonumber
    \eea        

    Because $\rho$ is unitary and $\Tr \sigma = 1$, $\EE_{x \in G} \| \rho(x) \|_\sigma^2 = 1$.
    For the second term, recall from \equref{eq:isometry_average} that the compression is the group average $V^* \rho^\prime(x) V = \EE_{y \in G} \rho^*(y)\rho(yx)$. Each summand $\rho^*(y)\rho(yx)$ is a product of unitaries, hence unitary, so $\|\rho^*(y)\rho(yx)\|_\sigma = 1$. By convexity of the $\sigma$-norm---the triangle inequality applied to the average---we obtain
    \be \label{eq:isometry_unitary_bound}
    \|V^* \rho^\prime(x) V\|_\sigma = \left\| \EE_{y \in G} \rho^*(y)\rho(yx) \right\|_\sigma \le \EE_{y \in G} \|\rho^*(y)\rho(yx)\|_\sigma = 1,
    \ee
    hence $\EE_{x \in G}\|V^* \rho^\prime(x) V\|_\sigma^2 \le 1$.

    Therefore using the bounds in Eq.~\ref{eq:isometry_unitary_bound} and Eq.~\ref{eq:inner_product_bound} by applying Lemma~\ref{gowers-hatami-correlation}, we have
    \be
        \EE_{x \in G}\left\|\rho(x)-V^{*} \rho^\prime(x) V\right\|_{\sigma}^{2} \le  2 - 2(1 - \varepsilon) = 2\varepsilon.
    \ee
\end{proof}
